\begin{center}
\refstepcounter{table}\label{tab:selection_function}
\begin{minipage}{\columnwidth}
\small\textbf{Table \thetable.} Gaia DR3 RVS selection-function
diagnostics for the catalogue.  The catalogue is selected from the
observed Gaia DR3 6D/RVS subset and is therefore not a volume-complete
Milky Way census.  $N_{\rm valid}$ is the number of sources for which
the GaiaUnlimited DR3-RVS implementation
\citep{CastroGinard2023} returned a finite selection-function weight at the source's $(G, BP{-}RP, l, b)$ cell;
the $N-N_{\rm valid}$ deficit is dominated by cells with zero parent
count.  $f_{n<10}$ is the fraction of $N_{\rm valid}$ sources whose
underlying GaiaUnlimited cell carries a parent count below 10 (the
regime in which the Beta$(k{+}1,\,n{-}k{+}1)$ posterior remains
competitive with its Beta$(1,1)$ prior: posterior standard error
$\gtrsim 0.16$ for $n<10$); $f_{\sum,n<10}$ is the share of
$\sum 1/S_{\rm RVS}$ contributed by those prior-dominated cells.  The
lower block restricts each tier to the inner-Galaxy quadrant
$l \in [330\degr,\,30\degr]$ where the catalogue concentrates
(Section~\ref{sec:sky}).  Inverse-selection weighted sums are reported
only as contextual observability diagnostics and are not used as
headline population counts.
\end{minipage}
\vspace{0.4ex}

\footnotesize
\setlength{\tabcolsep}{1pt}
\begin{tabular}{@{}lrrrrrr@{}}
\hline\hline
Sample & $N$ & $N_{\rm val}$ & $\Sigma 1/S$ & $\tilde S$ & $f_{<10}$ & $f_{\Sigma,<10}$ \\
\hline
\multicolumn{7}{l}{\textit{Full sample}} \\
Candidate pool          & 20{,}829 & 19{,}999 & 37{,}256.3 & 0.50 & --      & --      \\
Tier A                  & 276      & 274      & 435.7      & 0.67 & 80.3\%  & 86.7\%  \\
Tier B                  & 241      & 237      & 408.4      & 0.50 & 87.8\%  & 91.6\%  \\
Tier C                  & 1{,}318  & 1{,}261  & 2{,}324.0  & 0.50 & 86.3\%  & 89.5\%  \\
Tier A+B                & 517      & 511      & 844.0      & 0.67 & 83.8\%  & 89.0\%  \\
Tier A+B+C              & 1{,}835  & 1{,}772  & 3{,}168.1  & 0.50 & 85.6\%  & 89.4\%  \\
\hline
\multicolumn{7}{l}{\textit{Inner-Galaxy quadrant, $l\in[330\degr,\,30\degr]$}} \\
Tier A (inner)          & 35       & 34       & 47.0       & 0.67 & 58.8\%  & 66.2\%  \\
Tier A+B (inner)        & 79       & 78       & 122.6      & 0.67 & 66.7\%  & 73.8\%  \\
Tier A+B+C (inner)      & 715      & 664      & 1{,}260.1  & 0.50 & 82.7\%  & 85.2\%  \\
\hline
\end{tabular}
\end{center}
